\chapter{Conclusiones}\label{cap:conclusiones}

\section{Introducción}

Este capítulo recoge las conclusiones obtenidas tras el desarrollo del sistema \textbf{Sevillaneando}. Se revisan los objetivos planteados al inicio del proyecto y el grado en que han sido alcanzados, se analizan las lecciones técnicas aprendidas durante el desarrollo, se cuantifican y justifican las desviaciones respecto a la planificación original y se proponen líneas de trabajo futuro que permitirían ampliar y mejorar el sistema.

\section{Objetivos alcanzados}

El capítulo de estudio previo definió dos categorías de objetivos: los objetivos del proyecto, orientados al aprendizaje y la formación de la estudiante, y los objetivos del sistema, que describen las funcionalidades de la aplicación. A continuación se revisa el grado de cumplimiento de ambos.

\subsection{Objetivos del proyecto}

Los objetivos del proyecto hacen referencia a las metas formativas y profesionales perseguidas durante el desarrollo del TFG:

\begin{itemize}
    \item \textbf{Aplicación de conocimientos del Grado}: alcanzado. El proyecto ha integrado de forma transversal contenidos de diseño de software, arquitecturas de sistemas, ingeniería de requisitos y gestión de proyectos, aplicándolos en un contexto real y complejo.

    \item \textbf{Planificación y autogestión}: alcanzado. La gestión unipersonal del proyecto mediante metodología ágil adaptada ha requerido priorizar el backlog, revisar el alcance en función del tiempo disponible y tomar decisiones autónomas en cada sprint.

    \item \textbf{Experiencia en tecnologías modernas}: alcanzado. Se ha trabajado con React Native, NestJS, PostgreSQL/PostGIS y Firebase en un proyecto de extremo a extremo con despliegue real en producción.

    \item \textbf{Primer contacto con desarrollo móvil}: alcanzado. Se ha afrontado el ciclo de vida de una aplicación móvil, la adaptación a iOS y Android, y la compilación y distribución mediante Expo.

    \item \textbf{Consolidación de TypeScript}: alcanzado. TypeScript se ha utilizado de forma rigurosa tanto en el frontend como en el backend, profundizando en el tipado estático y las buenas prácticas del lenguaje.

    \item \textbf{Resolución de retos técnicos no triviales}: alcanzado. El sistema integra comunicación bidireccional en tiempo real, consultas geoespaciales con PostGIS, scraping automatizado de fuentes externas y un motor de recomendaciones multifactorial.

    \item \textbf{Documentación técnica de calidad}: alcanzado. El proyecto incluye memoria, manual de usuario, manual de desarrollo y manual de despliegue como parte de las entregas.
\end{itemize}

\subsection{Objetivos del sistema}

El objetivo principal del sistema era desarrollar una aplicación móvil multiplataforma que centralizase la oferta de ocio y cultura de Sevilla, con especial énfasis en dar visibilidad a eventos emergentes. A continuación se revisa el grado de cumplimiento de cada uno de los objetivos específicos:

\begin{itemize}
    \item \textbf{Listado y filtrado de eventos}: implementado completamente. La pantalla principal ofrece filtrado por categoría, precio, fecha y radio de distancia, con soporte para geolocalización en tiempo real.

    \item \textbf{Creación de eventos públicos y privados con moderación}: implementado completamente. Los eventos públicos siguen un flujo de moderación con estados \textit{Pendiente}, \textit{Aprobado} y \textit{Rechazado}. Los moderadores pueden revisar eventos pendientes en su listado propio y editar o eliminar todos los eventos públicos desde un listado específico, mientras que los eventos privados se aprueban automáticamente y son accesibles mediante enlace de invitación.

    \item \textbf{Edición y eliminación de eventos propios}: implementado completamente, incluyendo la lógica de vuelta a estado \textit{Pendiente} al editar un evento público aprobado y la posibilidad de eliminar eventos desde la pantalla de edición.

    \item \textbf{Geolocalización y mapa interactivo}: implementado completamente mediante OpenStreetMap y \texttt{react-native-maps}, con filtrado por radio y visualización de marcadores en mapa.

    \item \textbf{Interacción social}: implementado completamente. El sistema incluye chat de evento (sala grupal por evento), mensajería privada entre usuarios, valoraciones con puntuación y comentario, y módulo social basado en seguimiento, con listados diferenciados de seguidores, seguidos y amistades mutuas.

    \item \textbf{Sistema de recomendaciones personalizado}: implementado completamente. El motor aplica un algoritmo multifactorial ponderado que tiene en cuenta categorías de interés, historial de interacciones (guardar, asistir, compartir, ver, valorar), proximidad geográfica y temporal.

    \item \textbf{Rutas de eventos recomendadas}: implementado completamente. El motor genera itinerarios automáticos con tres estrategias: \textit{balanced}, \textit{walkable} y \textit{score}.

    \item \textbf{Creación de rutas propias}: implementado completamente. Los usuarios pueden crear rutas visibles para el resto de la comunidad, que pueden ser valoradas por otros usuarios.

    \item \textbf{Notificaciones}: implementado completamente. El sistema envía notificaciones por aprobación o rechazo de eventos, eventos cercanos, eventos próximos, nuevos mensajes privados, nuevos seguidores y recomendaciones generadas. Además, las notificaciones relacionadas con perfiles o eventos incluyen acciones para navegar directamente al perfil o al evento correspondiente.

    \item \textbf{Documentación automática de la API con Swagger}: implementado completamente. El backend expone un contrato OpenAPI completo mediante \texttt{@nestjs/swagger}. La documentación interactiva se genera automáticamente a partir de los decoradores \texttt{@ApiTags}, \texttt{@ApiOperation}, \texttt{@ApiParam}, \texttt{@ApiBearerAuth} y \texttt{@ApiProperty} añadidos a los diez controladores y los DTOs principales, y queda accesible en la ruta \texttt{/api/docs} de la API.
\end{itemize}

En cuanto a los objetivos del sistema no alcanzados, la generación automática de carteles con IA y la moderación de imágenes mediante NSFWJS (detección de contenido inapropiado) quedaron fuera del alcance final del proyecto, tal como se detalló en el análisis de desviaciones del capítulo de estudio previo.

\section{Resultados obtenidos}

El resultado principal del proyecto es una aplicación móvil multiplataforma funcional para el descubrimiento, creación y gestión de eventos en Sevilla. La solución final integra una aplicación cliente desarrollada con React Native y Expo, un backend basado en NestJS, una base de datos PostgreSQL con extensión PostGIS y distintos servicios externos para autenticación, almacenamiento de imágenes, mapas, geocodificación y obtención de eventos desde fuentes externas.

Desde el punto de vista funcional, el sistema permite consultar eventos, filtrarlos por criterios relevantes, visualizarlos en un mapa, crear, editar y eliminar eventos públicos y privados, moderar contenido, participar en chats, enviar mensajes privados, buscar usuarios, consultar sus perfiles, seguir o dejar de seguir a otras personas, revisar seguidores, seguidos y amistades mutuas, ver los eventos a los que asisten, valorar eventos, recibir notificaciones accionables y obtener recomendaciones personalizadas. Además, se han incorporado rutas de eventos recomendadas y rutas creadas por los propios usuarios, ampliando el alcance de la aplicación más allá de un listado convencional de actividades.

Desde el punto de vista documental, el proyecto ha generado una memoria completa, un manual de usuario, un manual de desarrollo, un manual de despliegue, actas de reunión, informes de avance y cambio, el reporte de horas de Clockify y la declaración de autoría correspondiente. Todo ello permite presentar el trabajo no solo como una aplicación terminada, sino como un proyecto de ingeniería documentado, planificado, probado y preparado para su mantenimiento o evolución futura.

\section{Lecciones aprendidas}

Los objetivos del proyecto definidos al inicio del TFG establecían compromisos de aprendizaje concretos. Esta sección reflexiona sobre qué ha aportado realmente cada uno de ellos.

\begin{itemize}
    \item \textbf{Aplicar en un contexto real los conocimientos del Grado.} El proyecto ha demostrado que las asignaturas de diseño de software, arquitecturas, ingeniería de requisitos y gestión de proyectos no son compartimentos estancos. Aplicarlas simultáneamente en un contexto con consecuencias reales ---una mala decisión de diseño en el modelo de datos repercutía directamente en las semanas siguientes--- ha consolidado esos conocimientos de una forma que ningún ejercicio académico aislado puede reproducir.

    \item \textbf{Desarrollar la capacidad de planificación, seguimiento y autogestión.} Gestionar el proyecto sin equipo ha evidenciado que la priorización del backlog es una habilidad no trivial. La combinación de estimación por tallas de camiseta y seguimiento real con Clockify ha sido fundamental para detectar desviaciones a tiempo y replantear el alcance de forma razonada, en lugar de acumular deuda y tomar decisiones precipitadas al final.

    \item \textbf{Adquirir experiencia práctica en tecnologías modernas.} Trabajar con React Native, NestJS, PostgreSQL/PostGIS y Firebase en un proyecto de extremo a extremo ha confirmado que la curva de aprendizaje real de estas tecnologías está en su integración conjunta, no en cada una por separado. Las fricciones aparecen en los bordes: entre el frontend y el backend, entre la autenticación y los guards, entre el ORM y las funciones espaciales.

    \item \textbf{Adentrarse por primera vez en el desarrollo de aplicaciones móviles.} El ciclo de vida de una aplicación móvil, las diferencias de comportamiento entre iOS y Android, el sistema de permisos y el proceso de compilación con Expo son conceptos sin equivalente directo en el desarrollo web. La lección más relevante ha sido aprender a leer documentación de alto nivel rápidamente y distinguir qué hay que entender a fondo de qué puede tratarse inicialmente como caja negra.

    \item \textbf{Consolidar el uso de TypeScript en un proyecto complejo de extremo a extremo.} Usar TypeScript de forma rigurosa en frontend y backend a la vez ha sido cualitativamente distinto a su uso en ejercicios académicos. Los DTOs con validación, los guards tipados de NestJS y los tipos de retorno de los hooks de React han hecho que muchos errores de integración fueran detectables en tiempo de compilación. El valor real de TypeScript no está en los tipos básicos, sino en modelar contratos entre módulos y hacerlos verificables automáticamente.

    \item \textbf{Enfrentar y resolver retos técnicos no triviales.} Cada módulo complejo ---scraping, recomendaciones, rutas, chat en tiempo real, consultas geoespaciales--- ha requerido una fase de investigación previa no contemplada en la estimación inicial. La incertidumbre técnica es un riesgo de primer orden en ingeniería del software: prototipar y validar un enfoque antes de comprometerse con él en el backlog es una práctica que este proyecto ha confirmado como imprescindible.

    \item \textbf{Producir documentación técnica de calidad.} Redactar la memoria, los manuales y los anexos ha supuesto una parte significativa del esfuerzo total. Esta experiencia ha cambiado la percepción de la documentación: no es un trámite que se hace al final, sino una actividad que revela inconsistencias en el diseño y obliga a articular decisiones que de otro modo quedarían implícitas en el código.
\end{itemize}

\section{Desviaciones respecto a la planificación}

Las desviaciones producidas durante el proyecto ---incorporación de CI/CD, mayor complejidad del scraping, no implementación de NSFWJS y generación de carteles con IA, y corrección de compatibilidad con Android--- se describen en detalle, con su impacto cuantitativo y justificación, en la sección~\ref{sec:desviaciones} del capítulo~\ref{cap:planificación}.

En conjunto, el proyecto se completó dentro del plazo establecido y con el alcance funcional prioritario cubierto al 100\%.

Con la perspectiva que da el cierre del proyecto, estas desviaciones revelan tres lecciones de fondo sobre la naturaleza de los proyectos de software:

\begin{itemize}
    \item \textbf{La incertidumbre técnica es inherente, no un fallo de planificación.} Las desviaciones más grandes ---el scraping y la compatibilidad con Android--- no eran predecibles con la información disponible al inicio. En ingeniería del software, una parte del esfuerzo real solo se descubre durante la implementación. Aceptar esto como una característica del trabajo, y no como un error de estimación, es una forma de madurez profesional que este proyecto ha ayudado a interiorizar.

    \item \textbf{Las dependencias ocultas entre módulos son más peligrosas que la complejidad individual.} El impacto del scraping sobre el sistema de rutas no era visible en el diseño inicial porque cada módulo parecía independiente. La lección no es solo técnica ---analizar dependencias antes de estimar--- sino de gestión: cuando dos módulos comparten datos o supuestos, un retraso en uno se convierte automáticamente en un retraso en el otro, y eso hay que anticiparlo al planificar el orden de desarrollo.

    \item \textbf{Entregar algo que funciona vale más que cumplir el plan original.} La decisión de no implementar las funcionalidades de la Fase 5 no fue un fracaso: fue una elección consciente de priorizar la calidad del núcleo sobre la cantidad de características. Esta capacidad de ajustar el alcance sin perder el objetivo principal es, probablemente, la lección de gestión más transferible de todo el proyecto.
\end{itemize}

\section{Consideraciones de seguridad y privacidad}

Durante el desarrollo se han aplicado medidas de seguridad a distintos niveles del sistema:

\begin{itemize}
    \item \textbf{Autenticación}: gestionada íntegramente por Firebase Authentication, que garantiza el almacenamiento seguro de credenciales y la emisión de tokens JWT de corta duración.
    \item \textbf{Autorización}: todos los endpoints protegidos verifican el token en cada petición mediante un guard de Firebase. Los endpoints de moderación y administración aplican un guard adicional de roles.
    \item \textbf{Limitación de velocidad}: se aplica \textit{throttling} en tres niveles (estricto, moderado y subida) mediante \texttt{@nestjs/throttler} para prevenir abusos y ataques de fuerza bruta.
    \item \textbf{Validación de entrada}: los DTOs del backend utilizan \texttt{class-validator} para validar y sanear todos los datos recibidos del cliente antes de procesarlos.
    \item \textbf{Privacidad de datos}: la aplicación no almacena contraseñas (delegadas en Firebase) ni datos de pago. Los datos de ubicación del usuario solo se utilizan para consultas espaciales y no se registran con seguimiento histórico.
    \item \textbf{Eventos privados}: el acceso a eventos privados requiere un enlace de invitación único generado en el backend, que no es predecible ni enumerable.
    \item \textbf{Row-Level Security en Supabase}: tras recibir un aviso de seguridad crítico de Supabase, se activó la política de \textit{Row-Level Security} (RLS) en las 18 tablas del esquema público de la base de datos. Sin esta medida, cualquier persona con la URL del proyecto y la clave anónima podía leer, modificar y eliminar datos directamente a través de la API REST de Supabase, sin pasar por el backend. El backend NestJS no se ve afectado, ya que se conecta como \textit{service role}, que omite las políticas RLS por diseño.
\end{itemize}

\subsection{Limitaciones respecto al RGPD}

La aplicación recoge y almacena datos personales sensibles sujetos al Reglamento General de Protección de Datos (RGPD, Reglamento UE 2016/679): ubicación geográfica en tiempo real, fotografías de perfil, historial de interacciones con eventos y mensajes privados entre usuarios. En el contexto de este TFG, la aplicación no está publicada en producción para el público general, por lo que estas obligaciones legales no son exigibles en esta fase. No obstante, se identifican a continuación las medidas implementadas y las limitaciones conocidas que deberían abordarse antes de cualquier despliegue real.

\textbf{Medidas implementadas.} La pantalla de registro incluye un checkbox de aceptación de los términos de uso y la política de privacidad, basando el tratamiento de datos en el consentimiento del usuario (RGPD art. 6.1.a). La fecha y hora de dicha aceptación se registran en la base de datos (\texttt{privacyAcceptedAt}) desde el primer inicio de sesión, lo que permite acreditar el consentimiento si fuera necesario. El endpoint de eliminación de cuenta (\texttt{DELETE /users/me/firebase}) propaga la baja a Firebase Authentication y, desde la fase de cierre del proyecto, también elimina la fotografía de perfil almacenada en Cloudinary mediante el \texttt{publicId} persistido en la entidad de usuario. Además, el backend incorpora un \texttt{PurgaDatosScheduler} expuesto mediante \texttt{POST /scheduler/run-purga-datos}; en el despliegue se configura un cron externo diario (\texttt{0 4 * * *}) que invoca dicho endpoint para eliminar mensajes privados huérfanos —aquellos cuyo emisor y receptor han sido eliminados— con más de 30 días de antigüedad.

\textbf{Limitaciones conocidas:}

\begin{itemize}
    \item \textbf{Consentimiento explícito para la geolocalización}: la aplicación solicita el permiso de ubicación del sistema operativo (iOS/Android), pero no incluye una pantalla de consentimiento informado que explique al usuario para qué se utiliza su localización, durante cuánto tiempo y con qué terceros podría compartirse. El RGPD exige que el consentimiento sea libre, específico, informado e inequívoco.

    \item \textbf{Política de retención de datos}: el sistema implementa una purga diaria de mensajes privados huérfanos (cuando ambas cuentas, emisor y receptor, han sido eliminadas) con más de 30 días de antigüedad, ejecutada mediante un disparador externo para no depender de que el proceso del backend permanezca activo de forma continua. Sin embargo, los mensajes privados en los que solo una de las partes ha eliminado su cuenta quedan huérfanos parcialmente (el campo del usuario eliminado se pone a \texttt{NULL} por la política \texttt{SET NULL}), pero el contenido del mensaje no se elimina. Una política de retención completa debería especificar que el contenido se anonimiza también cuando únicamente una de las partes ha eliminado su cuenta, o que los mensajes se eliminan en cascada al eliminar cualquiera de los dos usuarios.

    \item \textbf{Procedimiento formal de ejercicio del derecho al olvido}: aunque la eliminación técnica de la cuenta y los datos principales está implementada, no existe un canal documentado para que el usuario pueda ejercer formalmente este derecho frente a la responsable del tratamiento, ni un registro de solicitudes recibidas y atendidas, tal como exige el RGPD.
\end{itemize}

Estas limitaciones se documentan como trabajo futuro en la sección siguiente y constituyen una de las tareas prioritarias antes de un lanzamiento real de la aplicación.


\section{Trabajo futuro}

A lo largo del desarrollo han surgido ideas de mejora y nuevas funcionalidades que, por restricciones de tiempo o alcance, no han podido implementarse en esta versión del sistema. Se organizan a continuación en dos bloques: mejoras sobre funcionalidades existentes y nuevas funcionalidades.

\subsection{Mejoras sobre funcionalidades existentes}

\begin{itemize}
    \item \textbf{Cumplimiento RGPD}: las medidas básicas de consentimiento (checkbox en el registro con \texttt{privacyAcceptedAt}), derecho al olvido para imágenes de perfil (eliminación en Cloudinary al borrar la cuenta) y la purga diaria de mensajes privados huérfanos mediante cron externo (cuando ambos usuarios han sido eliminados) ya están implementadas. Las tareas pendientes antes de un despliegue real son: implementar una pantalla de consentimiento informado específica para la geolocalización, completar la política de retención para cubrir también los mensajes en los que solo una de las partes ha eliminado su cuenta, y establecer un procedimiento formal y documentado para el ejercicio del derecho al olvido. Estas limitaciones se identifican en la sección de consideraciones de seguridad y privacidad.

    \item \textbf{Moderación de contenido inapropiado (NSFW)}: Completar esta funcionalidad evitaría la subida de imágenes inapropiadas en eventos y chats, mejorando la seguridad del contenido sin requerir revisión manual para cada imagen.

    \item \textbf{Verificación de correos electrónicos}: actualmente, el registro de usuarios no incluye un paso de validación del correo mediante enlace de confirmación. Implementar esta funcionalidad mediante Firebase Auth o un servicio de correo transaccional (SendGrid, Mailgun) aumentaría la confiabilidad de las direcciones de correo registradas y permitiría re-enviar invitaciones o notificaciones críticas con mayor seguridad. Además, facilitaría la recuperación de cuentas y reduciría registros fraudulentos o con errores tipográficos.

    \item \textbf{Cobertura de pruebas}: el plan de pruebas actual cubre los flujos principales pero deja fuera casos límite en el motor de recomendaciones, el scraping y la lógica de moderación. Ampliar la suite con pruebas de carga y pruebas de regresión automatizadas aumentaría la fiabilidad del sistema en producción.
\end{itemize}

\subsection{Nuevas funcionalidades}

\begin{itemize}
    \item \textbf{Generación de carteles con IA}: identificada desde el inicio como objetivo de la Fase 5, esta funcionalidad permitiría a los usuarios generar automáticamente un cartel visual para su evento a partir del título, la descripción y las imágenes proporcionadas, integrando un modelo generativo de imagen a través de una API externa.

    \item \textbf{Modo sin conexión básico}: actualmente no existe detección de red ni caché local. Mediante \texttt{AsyncStorage} o React Query se podrían almacenar localmente los últimos eventos cargados y mostrarlos sin conexión, mejorando la experiencia de usuario en entornos con conectividad intermitente.

    \item \textbf{Achievements y gamificación}: dado que el sistema ya registra visitas, asistencias, guardados y compartidos, sería natural introducir un sistema de insignias: \textit{Explorador} (10 eventos distintos visitados), \textit{Foodie} (5 eventos de gastronomía asistidos), \textit{Organizador} (3 eventos propios aprobados), entre otras. Este mecanismo aumentaría la retención y la participación activa de los usuarios.

    \item \textbf{Módulo avanzado de amigos}: actualmente el sistema dispone de un módulo social basado en seguimiento, con búsqueda de usuarios, perfiles públicos, opción de seguir o dejar de seguir, distinción entre seguidores, seguidos y amistades mutuas, mensajería privada y consulta de eventos a los que asisten otros usuarios. Como ampliación futura, se podría incorporar una gestión más explícita de contactos cercanos, invitaciones directas a eventos y recomendaciones sociales basadas en intereses compartidos.

    \item \textbf{Integración con el calendario nativo del dispositivo}: la pantalla \texttt{CalendarEventsScreen} ya existe pero no hay integración con el calendario nativo del teléfono. Mediante \texttt{expo-calendar} sería posible añadir directamente los eventos a los que el usuario se apunta en Google Calendar o Apple Calendar, sin salir de la aplicación.

    \item \textbf{Reseñas con fotos}: las reseñas (\texttt{Resena}) actualmente contienen puntuación y comentario de texto. Permitir adjuntar fotografías en las valoraciones aumentaría la credibilidad de las reseñas y la confianza de los usuarios a la hora de apuntarse a un evento.

    \item \textbf{Asistencia mediante código QR}: integrar un sistema de códigos QR que el organizador pueda mostrar en el evento permitiría registrar la asistencia de forma sencilla. Combinado con un enlace de inscripción externo (Eventbrite, formularios propios), facilitaría la interoperabilidad con otras plataformas de gestión de eventos.

    \item \textbf{Cuentas de empresa y promoción de eventos}: actualmente todos los usuarios tienen el mismo tipo de cuenta. Introducir un perfil de empresa o entidad permitiría a organizaciones y negocios registrarse con un perfil verificado y destacado, con la posibilidad de promocionar sus eventos dentro de la plataforma mediante un modelo de visibilidad preferente.

    \item \textbf{Soporte de eventos online}: el modelo de evento actual está orientado a actividades presenciales, con dirección física y ubicación geográfica como atributos centrales. Ampliar el modelo para contemplar eventos online (conciertos en streaming, talleres por videoconferencia, presentaciones virtuales) requeriría añadir un campo de tipo de celebración (\textit{presencial}, \textit{online} o \textit{híbrido}), sustituir la ubicación por un enlace de acceso y adaptar el motor de filtrado para que los eventos online sean visibles independientemente del radio de búsqueda del usuario. Esta extensión ampliaría significativamente el catálogo disponible y abriría la plataforma a organizadores que operan sin espacio físico fijo.

    \item \textbf{Pasarela de pago integrada}: actualmente la aplicación muestra el precio de un evento como información orientativa pero no gestiona la compra de entradas. Integrar una pasarela de pago (por ejemplo, Stripe) permitiría al usuario adquirir su entrada directamente desde la pantalla de detalle del evento, sin necesidad de redirigirle a una plataforma externa. Esta funcionalidad requeriría ampliar el modelo de datos con entidades de pedido y ticket, implementar los webhooks de confirmación de pago, y gestionar la emisión de entradas en formato QR escaneables en el acceso al evento. Desde el punto de vista del organizador, ofrecería un canal de venta propio con control sobre la lista de asistentes.

    \item \textbf{Soporte multiidioma (i18n)}: la aplicación está actualmente disponible únicamente en español. Añadir internacionalización mediante una librería como \texttt{react-i18next} permitiría ofrecer la interfaz en varios idiomas (inglés, francés, etc.), ampliando el alcance de la plataforma a turistas y visitantes internacionales, un perfil de usuario especialmente relevante dado el carácter turístico de Sevilla. La implementación requeriría centralizar todos los literales de la interfaz en archivos de traducción por idioma y añadir un selector de idioma en la configuración del perfil.
\end{itemize}

\section{Uso de herramientas de inteligencia artificial}

Durante el desarrollo de este proyecto se han utilizado herramientas de inteligencia artificial como apoyo en determinadas tareas. A continuación se detalla el uso realizado, las herramientas empleadas y el proceso de verificación seguido en cada caso, en cumplimiento con los principios de transparencia exigidos por la Escuela Técnica Superior de Ingeniería Informática de la Universidad de Sevilla.

\subsection{Herramientas utilizadas}

\begin{itemize}
    \item \textbf{GitHub Copilot}: asistente de autocompletado de código integrado en el entorno de desarrollo. Se ha utilizado para sugerir fragmentos de código repetitivo (DTOs, validadores, controladores NestJS con estructura similar) y para acelerar la escritura de código boilerplate en el frontend (componentes con estructura análoga a otros ya existentes). Además, se ha empleado puntualmente para proponer mejoras visuales en componentes de interfaz, como ajustes de estilos y disposición de elementos.

    \item \textbf{Claude (Anthropic)}: modelo de lenguaje utilizado para resolver dudas técnicas concretas (p.ej., uso de funciones espaciales de PostGIS, configuración de guardias de autenticación en NestJS), proponer mejoras visuales en la interfaz y como guía en procesos técnicos más elaborados, como el despliegue en producción y la revisión de vulnerabilidades de seguridad. Asimismo, la redacción de la memoria ha sido revisada con ayuda de esta herramienta para mejorar la expresión y coherencia del texto. De forma específica, se utilizó para generar el contenido inicial de la tabla de endpoints de la API REST (Tabla~\ref{tab:api}) a partir de la estructura de controladores y DTOs del backend, dado el carácter mecánico y repetitivo de esta tarea; el resultado fue revisado, corregido y validado íntegramente por la autora.
\end{itemize}

\subsection{Alcance del uso}

El uso de estas herramientas ha sido complementario y supervisado en todo momento. No se ha delegado en ellas la toma de decisiones de diseño, arquitectónicas ni funcionales; éstas han sido adoptadas íntegramente por la autora. El código generado automáticamente ha sido siempre revisado, comprendido y adaptado antes de integrarse en el proyecto, nunca copiado directamente sin verificación. Las guías proporcionadas para procesos como el despliegue o la revisión de seguridad han servido como referencia técnica, pero las decisiones finales y su implementación han sido responsabilidad de la autora.

La redacción de la memoria ha sido elaborada por la autora. Las herramientas de IA se han utilizado como apoyo para mejorar la expresión y coherencia del texto, pero el contenido, la estructura y las conclusiones son propios.

\subsection{Verificación}

Todo el código sugerido por herramientas de IA ha pasado por el mismo proceso de revisión que el código escrito manualmente: revisión de lógica, ejecución de pruebas unitarias y de integración, y comprobación del comportamiento en la aplicación. No se ha aceptado ningún fragmento cuyo funcionamiento no haya podido ser explicado y verificado por la autora.

\section{Conclusión final}

\textit{Sevillaneando} es el resultado de un proceso de ingeniería completo: desde la identificación de un problema real ---la fragmentación de la información sobre eventos en Sevilla--- hasta la entrega de un sistema funcional, desplegado y documentado. El proyecto ha cubierto todos los objetivos prioritarios definidos al inicio, integrando tecnologías actuales de desarrollo móvil, backend, geolocalización, comunicación en tiempo real y recomendaciones personalizadas.

Más allá del resultado técnico, el desarrollo ha supuesto un aprendizaje profundo en la gestión de proyectos software en solitario: tomar decisiones de alcance bajo incertidumbre, priorizar con criterio y entregar algo que funciona sobre algo que es perfecto en papel. Las funcionalidades no implementadas no son fracasos sino decisiones conscientes de gestión del alcance, y quedan documentadas como trabajo futuro para una versión posterior del sistema.
